\section{The problem: sound event detection}

Sound event detection (SED) takes an audio clip and returns a set of events
\begin{equation}
\mathcal{E} = \{(c_j,\; t^{\mathrm{start}}_j,\; t^{\mathrm{end}}_j)\},
\end{equation}
where $c_j$ is the event class and the two timestamps are its boundaries.

There are two common designs. \emph{Event-based} systems directly predict a set
of events, but need complex objectives such as Hungarian matching.
\emph{Frame-wise} systems predict, for every class $c$ and frame $t$, a
probability that the event is active:
\begin{equation}
\hat p^{\pres}_{c,t} = P(e_{c,t}=1).
\end{equation}
Frame-wise models are the dominant paradigm because they train well with
multiple-instance learning and simple architectures such as a convolutional
recurrent neural network (CRNN), and they work with weakly labelled data.

Their weakness is converting frame probabilities into events. The standard
recipe is: (i) smooth the probability curve, usually with a median filter;
(ii) threshold it; (iii) group consecutive active frames into events. The filter
length and threshold are tuned on a validation set after training. This
post-processing is non-differentiable and can discard information the model
learned. The paper's goal is to remove that separate tuning step by learning
event boundaries and proposals inside the model.

